% =====================================================================
% Cover letter VARIANT B (reproducibility/rigor-led) — NMR in Biomedicine
% EIC: Professor John R. Griffiths
% Built with tectonic: `tectonic cover_letter_reproducibility.tex`
% =====================================================================
\documentclass[11pt]{letter}
\usepackage[margin=1in]{geometry}
\usepackage{microtype}
\usepackage[hidelinks]{hyperref}

\signature{Avery Karlin, BA \\ Department of Applied Physics and Applied Mathematics,
Columbia University \\ Department of Computer Science, University of Colorado Boulder \\
\texttt{ak5232@columbia.edu}}

\address{Avery Karlin \\ Columbia University \\ New York, NY, USA}

\begin{document}

\begin{letter}{Professor John R.\ Griffiths \\ Editor-in-Chief, \textit{NMR in Biomedicine}}

\opening{Dear Professor Griffiths,}

I am pleased to submit the manuscript ``Pseudo-diffusion is weakly identifiable in
practice: a reproducible boundary-railing signature of $D^{*}$ in open in-vivo IVIM, and
the shape-aware intervals that resolve it'' for consideration in \textit{NMR in
Biomedicine}.

I want to foreground how the work was done, because its method of construction is part of
its contribution. The study was built to a strict run-then-write discipline: no number
appears in the manuscript that was not first produced by a seeded run, and every reported
value is frozen in a tracked file and regenerable with a single command. The entire
quantitative core was re-derived in a clean-room reimplementation written independently
of any prior analysis code, on openly licensed in-vivo data (OSIPI TF2.4, Zenodo
14605039, SHA-256-verified; an independent TCGA-LIHC liver cohort via The Cancer Imaging
Archive). A reader can reproduce the lead result --- and every figure --- from scratch.

This discipline is also what makes the scientific claim trustworthy. The central result
is a real, reproducible phenomenon: a conventional least-squares IVIM fit rails the
pseudo-diffusion coefficient $D^{*}$ onto a parameter bound in 54.2\% (95\% CI
52.0--56.4) of high-SNR voxels of an open abdomen scan, reproducing a previously reported
54.7\% without using its code, and recurring at 44--73\% across independent abdominal and
liver cohorts. Crucially, the same rigour led us to \emph{retract} a previously reported
headline. An earlier pair of marginal $D^{*}$ coverage figures (0.30 / 0.67) did not
reproduce; we traced the discrepancy to an undocumented standard-deviation convention,
and we now report $D^{*}$ coverage \emph{conditionally}, per true-$D^{*}$ tercile, under
an explicitly stated honest Cram\'er--Rao convention. The honest finding --- a
conditional under-coverage concentrated in the high-$D^{*}$ regime (0.63), resolved by
shape-aware quantile intervals and an amortized posterior up to an irreducible
identifiability wall --- is less dramatic than the retracted number, and far more
defensible. We think that trade is exactly the right one for the literature.

The manuscript fits the scope of \textit{NMR in Biomedicine} in both its principal
categories: it is (a) an original research study of an in-vivo IVIM phenomenon and (b) a
methodological contribution to how quantitative-MR uncertainty should be measured and
honestly reported, with a worked, convention-explicit example. The reproducibility module
is offered to the community for benchmarking future IVIM uncertainty methods.

The manuscript is original, has not been published previously, and is not under
consideration elsewhere. The author declares no competing interests and takes full
responsibility for the integrity of the work; a declaration of AI-assisted writing is
included. I would be glad to suggest suitable reviewers and to answer any questions.

\closing{Sincerely,}

\end{letter}
\end{document}
